\documentclass[11pt,english]{article}
\bibliographystyle{jf}
\usepackage[longnamesfirst]{natbib}
%\usepackage[T1]{fontenc}
\usepackage{fontenc}
\usepackage[utf8]{inputenc}
\usepackage[bottom]{footmisc}
\usepackage{tabularx}
\usepackage{booktabs}
\usepackage{amsmath,enumerate}
\usepackage{amsthm}
\usepackage{dsfont}
\usepackage{amssymb}
\usepackage{mathtools}
\usepackage{mathrsfs}
\usepackage{graphicx}
\usepackage{natbib}
\usepackage[labelfont=bf,justification=justified,font=small,skip=2pt,labelsep=period]{caption}
%\usepackage[labelfont=bf,justification=justified,labelsep=period]{caption}
\usepackage{subcaption}
\usepackage{setspace}
\usepackage{arydshln}
\usepackage{chngpage}
\usepackage{float,appendix}
\usepackage{afterpage}
\usepackage{verbatim}
\usepackage{indentfirst}
\usepackage[margin=1in]{geometry}
\usepackage[usenames, dvipsnames]{color}
\usepackage{hyperref}
\usepackage{adjustbox}
% \usepackage{longtable}
\usepackage{booktabs, makecell, longtable,ltxtable}
\usepackage{lipsum}
\usepackage{pdflscape}
\usepackage{siunitx}
\usepackage{chngcntr}
\usepackage{lipsum}
\usepackage{tikz}
\usepackage{pdflscape}
\usepackage[normalem]{ulem}
\usepackage{chngcntr,apptools}
\AtAppendix{\counterwithin{lemma}{section}}
\usepackage{siunitx}
\usepackage{tabularx}
\sisetup{table-format=-1.3, table-space-text-post={**}}
\usepackage{xr}
\usepackage{etoolbox}
\usetikzlibrary{positioning}
\newcommand\fnsep{\textsuperscript{,}}
\usepackage{tikz}
\usepackage{pdflscape}
\usepackage[normalem]{ulem}
\usepackage{chngcntr,apptools}
\AtAppendix{\counterwithin{lemma}{section}}
\usepackage{siunitx}
\usepackage{tabularx}
\usepackage{xr}
\usepackage{etoolbox}
\usetikzlibrary{positioning}

%% UCLA colors: 
% https://brand.ucla.edu/wp-content/uploads/ucla-color-palette-v10.pdf
%\definecolor{UCLA}{RGB}{50, 132, 191}   % UCLA Blue primary
%\definecolor{UCLA}{RGB}{52, 123, 173}   % UCLA Blue text only
\definecolor{UCLAblue}{RGB}{30, 75, 135} % secondary color	
\definecolor{UCLAgold}{RGB}{255, 232, 0} % UCLA Gold

\definecolor{usccardinal}{rgb}{0.6, 0.0, 0.0}
\definecolor{Matlabblue}{rgb}{0,0.447,0.741}
\definecolor{bleudefrance}{rgb}{0.19,0.55, 0.91}
\definecolor{cobalt}{rgb}{0.0, 0.28,0.67}
\definecolor{darkblue}{rgb}{0.0, 0.0,0.55}
\definecolor{darkcerulean}{rgb}{0.03,0.27, 0.49}
\definecolor{darkpowderblue}{rgb}{0.0,0.2, 0.6}
\definecolor{bleudefrance}{rgb}{0.19,0.55, 0.91}
%\usepackage[colorlinks=true,urlcolor=blue,linkcolor=blue,citecolor=blue]{hyperref}

\hypersetup{
	colorlinks,
	linkcolor={UCLAblue},
	citecolor={UCLAblue},
	urlcolor={UCLAblue},
}

\bibpunct{(}{)}{;}{a}{,}{,}
\newcommand{\E}{\mathbb{E}}
\newcommand{\var}{\mathrm{Var}}
\newcommand{\cov}{\mathrm{Cov}}
\setcounter{MaxMatrixCols}{20}
\newtheorem{theorem}{Theorem}
\newtheorem{lemma}{Lemma}
\newtheorem{defn}{Definition}
\newtheorem{prop}{Proposition}
\newtheorem{cor}{Corollary}
\urlstyle{same}
\usepackage{pgfplots} 
\pgfplotsset{compat=newest} 
\pgfplotsset{plot coordinates/math parser=false} 
% puts figures and tables at the end of the document
%\usepackage[nomarkers,nolists]{endfloat}               
% puts ONLY figures at the end of the document
%\usepackage[nomarkers,nolists,figuresonly]{endfloat}  
% allows more than one figure per page for endfloat
%\renewcommand{\efloatseparator}{\mbox{}}              
\usepackage{rotating}
% allows sideways tables and figures to be placed at the end of the document
%\DeclareDelayedFloatFlavor{sidewaystable}{table}
%\DeclareDelayedFloatFlavor{sidewaysfigure}{figure}

% \usepackage[nomarkers,nolists]{endfloat} 
% \renewcommand{\efloatseparator}{\mbox{}}              % allows more than one figure per page for endfloat

%\usepackage{mathpazo}
%\usepackage{mathpple}
%%%%%%%%%%%%%%%%%%%%%%%%%%%%%%%%
%\usepackage{palatino}
%%%%%%%%%%%%%%%%%%%%%%%%%%%%%%%%
\pgfplotsset{compat=newest} 
\pgfplotsset{plot coordinates/math parser=false} 
\newlength\figureheight 
\newlength\figurewidth 
\usepackage{mathtools}
\usetikzlibrary{calc}


%%%%%%%%%%%%%%%%%%%%%%%%%%
% Notes customization %
%%%%%%%%%%%%%%%%%%%%%%%%%%
\definecolor{webblue}{rgb}{0,0,0.6}
\definecolor{webgreen}{rgb}{0,.6,0}
\definecolor{webbrown}{rgb}{.6,0,0}


%- Commenting
\let\oldmarginpar\marginpar
\renewcommand\marginpar[1]{\-\oldmarginpar[\raggedleft\tiny #1]%
{\raggedright\tiny #1}}
\newcommand{\mknote}[1]{\marginpar{\textcolor{webblue}{\tiny{Mahyar: #1}}}}
\newcommand{\dsnote}[1]{\marginpar{\textcolor{webgreen}{\tiny{Dejanir: #1}}}}
\usepackage[en-US]{datetime2}
\usepackage{mathtools}
 \usepackage{xcolor}

\newcommand{\comments}[1]
{\par {\bfseries \color{blue} #1 \par}}

\usepackage[draft]{todonotes}

\title{\textbf{A Model with Equity and Risky Debt Claims }}
\author{ Dejanir Silva\thanks{Daniels School of Business, Purdue University, e-mail: \href{dejanir@purdue.edu}{dejanir@purdue.edu}.}}
\date{September 2025}
%\date{\today\\
%	\begin{small}
%		First draft: May 5, 2018
%\end{small}}
\onehalfspacing
\begin{document}
	\maketitle
	
\section{A Lucas Economy with a Levered Claim}

Consider an economy where the output of the representative firm follows a geometric Brownian motion:
\begin{equation*}
    \frac{d Y_t}{Y_t} = \mu dt + \sigma d Z_t.
\end{equation*}

The firm issues two types claims to finance its operation, capturing different equity and debt instruments. 
The first instrument is a claim on the firm's dividends. 
To capture the fact that dividends are more volatile than consumption (see, e.g., \citealt{Abel1999LeverageDividends}; \citealt{Campbell1999AssetPrices}), we follow \citet{MenzlySantosVeronesi2004} and assume the dividend process of the form $D_t = s_t Y_t$, where the share $s_t$ follows the process:
\begin{equation}
    d s_t = \kappa_s (\overline{s} - s_t)\, dt + \sigma_s s_t (1-s_t)\, dZ_t,
\end{equation}
given $s_0 \in (0,1)$, where $\overline{s} \in (0,1)$ and $\kappa_s > 0$. This ensures that $s_t \in [0,1]$ at all times. 

Ito's lemma implies that dividends follow the process:
\begin{equation}
    \frac{d D_t}{D_t} = \mu_D(s_t)\, dt + \sigma_D(s_t)\, dZ_t,
\end{equation}
where
\begin{equation}
    \sigma_{D}(s_t) = \sigma + \sigma_s (1-s_t), \qquad \qquad 
    \mu_D(s_t) = \mu + \kappa_s \left( \frac{\overline{s}}{s_t}-1\right) + \sigma \sigma_s(1-s).
\end{equation}

This implies that $\sigma_{D}(s_t) > \sigma$ for $\sigma_s > 0$. 
Therefore, dividends are more volatile than consumption. 
Moreover, this volatility process captures a form of ``leverage effect," where volatility increases after a negative shock. 
Finally, this process ensures that the share process is stationary, so consumption and dividends are co-integrated. 

The second instrument is a claim on the remaining payoff, $B_t = (1-s_t) Y_t$, which we refer to as a \textit{corporate bond}. 
Notice that the corporate bond payoff will be smoother than aggregate consumption, capturing the idea this is a safer claim. 

We denote the price of a claim on output by $P_{Y,t}$, the price of a claim on dividends by $P_{D,t}$, and the price of the corporate bond by $P_{B,t}$. 
As $B_t+D_t = Y_t$, we obtain $P_{B,t} +P_{D,t}= P_{Y,t}$.




\newpage
\singlespacing
%\bibliographystyle{jf}
\phantomsection\addcontentsline{toc}{section}{\refname}
\clearpage
%{\small
\bibliography{references.bib}
\end{document}